\section*{Part 2}
\subsection*{Q1}
Q1.1: We have more than double as much disease samples as healthy samples. See Appendix Figure~\ref{fig:Q1_1_label_distribution} for the label distribution plot.

Q1.2: Healthy images typically show clear, well-demarcated lung fields, even exposure, visible diaphragmatic contours, and no large dense opacities. Pneumonia images often show local or diffuse areas of increased density (consolidations), irregular, patchy opacities, possibly reduced lung volume or pleural effusion; affected regions are frequently located in the lower/mid lung fields. Key visual differences: focal or multilobar opaque regions in pneumonia versus largely clear, homogeneous lung fields in healthy cases; contrast and texture differ (anatomical details are often less discernible in pneumonia due to overlapping opacities). Representative example images are shown in Appendix Figure~\ref{fig:Q1_2_examples_grid}.

Q1.3: One plausible source of bias is the presence of acquisition or preprocessing artifacts that correlate with the label, such as differences in age groups, scanner settings, image contrast, or visible markers/text overlays. This can make the classifier rely on dataset-specific shortcuts instead of true pathology.

Q1.4: The preprocessing used for further analysis is: Convert images to grayscale, resize to a fixed resolution (224×224), normalize pixel intensities (mean=0.5, std=0.5), and train with light augmentation (random rotation, random resized crop) on the training split only.

\subsection*{Q2}
Q2.1: Our CNN is intentionally compact so it trains quickly on the course setup while still being expressive enough for attribution methods. The architecture consists of four convolutional blocks (32, 64, 128, 256 channels) with batch normalization, ReLU activation, and max pooling, followed by adaptive average pooling and a binary classifier. The model is optimized with weighted binary cross-entropy loss to handle class imbalance.

Q2.2: For reporting the performance on a test set, we created a confusion matrix and classification report (see Figure~\ref{fig:Q2_2_confusion_matrix}). The model achieves a test accuracy of 80.29\%, with a precision of 95.25\% and recall of 72.05\% for pneumonia detection. The F1-score is 0.8204. The confusion matrix shows 220 true negatives and 14 false positives (specificity: 94.02\%), with 281 true positives and 109 false negatives. This indicates the model is conservative about predicting pneumonia, which keeps false positives low, but it misses a notable number of pneumonia cases. In practice, this means it is better at confirming pneumonia when it predicts it than at ruling it out reliably.

\subsection*{Q3}
Q3.1: Integrated Gradients visualizations for sample images (5 healthy, 5 pneumonia) are in Appendix Figure~\ref{fig:Q3_1_integrated_gradients_normal_vs_pneumonia}. The method computes attributions using a zero baseline and 32 interpolation steps.

Q3.2: Integrated Gradients often highlight regions within the lung fields that correspond to the most visible consolidations (e.g., patchy or segmental opaque areas). However, in some cases the maps produce fine-grained or edge-focused patterns that also emphasize boundaries or contrast differences and do not exclusively indicate pathological regions. See Figure~\ref{fig:Q3_1_integrated_gradients_normal_vs_pneumonia}.

Q3.3: Attributions from Integrated Gradients are partially consistent: in cases with clear consolidations we observe similar highlighted regions in the affected lung area across different patients. However, the maps vary between subjects and depend on image quality, positioning, and contrast; for weak or diffusely distributed findings the patterns are often inconsistent. See Figure~\ref{fig:Q3_1_integrated_gradients_normal_vs_pneumonia}.

Q3.4: The choice of baseline image strongly affects Integrated Gradients: a zero baseline (black) emphasizes contrast changes relative to the black reference, while a mean or blurred baseline can attenuate or redistribute attributions. In our experiments, the zero baseline produces sharp, high-contrast attribution maps that highlight local discriminative features effectively.

\subsection*{Q4}
Q4.1: Grad-CAM overlays are shown in Appendix Figure~\ref{fig:Q4_1_gradcam}. The method visualizes the class-wise activation maps from the final convolutional layer, projected back to the input image resolution.

Q4.2: Grad-CAM typically highlights regions inside lung fields in our examples, showing anatomically plausible focus areas. Both healthy and pneumonia samples reveal activation in appropriate respiratory regions, though with varying spatial patterns.

Q4.3: Stability across samples is variable; some samples show consistent lung-field activation patterns within each class, while others reveal inconsistent or scattered activations. This suggests that while Grad-CAM captures class-relevant features, its spatial focus can vary substantially across individual test samples, even within the same diagnosis.

Q4.4: Compared to Integrated Gradients, Grad-CAM is coarser and more spatially focused on high-level features in the final convolutional layer. IG provides pixel-level signed attributions with finer detail, whereas Grad-CAM produces smoothed, spatially localized heatmaps. IG is sensitive to the baseline choice, while Grad-CAM depends on the learned feature maps of the architecture. Both methods show complementary strengths: IG reveals fine-grained patterns, while Grad-CAM identifies high-level discriminative regions.

\subsection*{Q5}
Q5.1: We retrained the CNN on randomly permuted training labels by shuffling the labels of all training samples while preserving image-label correspondence. The model was trained for 5 epochs using the same architecture, hyperparameters, and loss function as the original model to establish a baseline for comparison.

Q5.2: \textbf{Grad-CAM: Mixed result (partial pass).} After analyzing the 10 randomization comparison figures, Grad-CAM shows variable behavior. Approximately 60\% of samples exhibit noticeable differences between original and random-label CAM maps (colored difference heatmaps), indicating the method passes the sanity check for those samples. However, 40\% of samples show near-zero differences (black or very dark heatmaps), indicating failure---Grad-CAM produces identical maps regardless of label changes for these samples, suggesting it detects low-level features (edges, texture) independent of the diagnosis label.

\textbf{Integrated Gradients: Similar mixed result (partial pass).} The IG randomization test shows a similar pattern. Approximately 50\%-60\% of samples show substantial differences in signed attributions between original and random models (seismic colormap differences visible), suggesting IG is label-sensitive for those cases. The remaining 40\%-50\% show minimal differences or black/near-zero absolute differences, indicating IG fails the sanity check for those samples and may be capturing general patterns uninformed by the task labels.

\textbf{Overall verdict: Both methods show partial failure.} Neither method robustly passes the randomization test across all samples, indicating both can act as generic feature detectors (e.g., edge detectors, texture analyzers) independent of the pneumonia diagnosis task.

Q5.3: The data randomization test reveals a critical insight: both Integrated Gradients and Grad-CAM exhibit mixed trustworthiness. The finding that 40\--50\% of samples show no change in attribution maps after label randomization is concerning and aligns with prior work by Simonyan et al. (``Sanity Checks for Saliency Maps''), which demonstrated that naive saliency methods can fail this basic consistency test.

\textbf{Interpretation:} For samples showing colored/seismic differences in the comparison figures, the attribution maps are label-sensitive and likely capture decision-relevant features. However, for samples producing black difference heatmaps, the methods reduce to low-level feature detectors (e.g., Grad-CAM highlighting edges, IG emphasizing contrast boundaries) regardless of the labels. This suggests that while both methods can provide useful insights, they should be interpreted cautiously and only trusted for individual samples where the original and random attribution maps clearly diverge.

\textbf{Clinical implications:} From a medical deployment perspective, these results suggest that attributions should be validated on a per-sample basis rather than assumed universally trustworthy. Clinicians should view attribution maps as suggestive evidence rather than definitive proof of model reasoning. Methods that pass the sanity check more robustly (e.g., attention-based models or inherently interpretable architectures) would be preferable for high-stakes medical applications.

\section*{Part 3}
\subsection*{Q1}
Q1.2: \textbf{Moderate consistency between Integrated Gradients and Grad-CAM for Part 2.} Both methods identify lung fields as decision-relevant regions, showing broad agreement at the global level. However, they differ substantially at the local level. IG reveals fine-grained, pixel-level patterns (edges, texture, consolidation boundaries), while Grad-CAM produces coarser, spatially localized heatmaps from high-level features. In healthy samples, both tend to show relatively distributed activation; in pneumonia samples, both focus on affected regions, but IG highlights finer consolidation details while Grad-CAM captures broader spatial patterns. This partial consistency suggests both methods identify genuine diagnostic features, but their architectural dependencies (IG: gradient-based, Grad-CAM: feature-map based) cause them to emphasize different levels of detail.

\subsection*{Q2}
Q2.2: \textbf{For pneumonia diagnosis (Part 2), I would convince a doctor to trust Grad-CAM over Integrated Gradients.} While Grad-CAM failed the full randomization sanity check, it offers superior clinical interpretability. When Grad-CAM produces colored, spatially focused heatmaps (not black difference maps), it clearly highlights suspected affected lung regions that a radiologist can visually verify. The localized activation is easier for a clinician to defend: ``The model attends to the lower-left lobe, which matches the radiological findings.'' By contrast, IG's fine-grained pixel-level attributions are harder to communicate and interpret clinically. To build trust, I would: (1) show Grad-CAM overlays alongside the original radiographs for cases where the randomization test passes (colored difference maps), (2) explicitly note which samples fail the sanity check (black differences), treating those as unreliable, and (3) emphasize that high-attention regions should complement, not replace, radiologist judgment. This transparency and acknowledgment of limitations would reassure a doctor that the model's reasoning is at least partially trustworthy and clinically aligned.

\subsection*{Q3}
Q3.2: \textbf{Yes, both IG and Grad-CAM attributions make strong intuitive and medical sense for pneumonia detection.} Both methods consistently highlight lung fields as decision-relevant, which is radiologically sound---pneumonia manifests as opacities in lung tissue. Both methods show stronger activation in pneumonia samples within the affected regions (lower/mid lobes, diffuse fields), matching classical radiological patterns. Consolidations, pleural effusions, and other pathological signatures naturally produce high attribution in both IG (contrast changes) and Grad-CAM (feature map activations). The methods do not, for example, highlight the cardiac silhouette or mediastinum as primary decision factors. This alignment between learned features and medical domain knowledge suggests the models have captured genuine diagnostic patterns rather than arbitrary statistical artifacts. The partial failure on the randomization test (40\--50\% of samples) indicates some generic feature detection, but the majority of evidence supports genuine pathology-based reasoning.

\subsection*{Q4}
Q4.2: \textbf{For practical pneumonia diagnosis deployment, I would choose Grad-CAM.} Although both Integrated Gradients and Grad-CAM show mixed sanity-check results, Grad-CAM is preferable for clinical deployment because: (1) \textbf{Interpretability}: Clinicians naturally understand spatial heatmaps; Grad-CAM outputs align with radiological reasoning (``what regions matter?''), making it easier to defend model decisions to regulatory bodies and patients. (2) \textbf{Computational efficiency}: Grad-CAM requires a single backward pass through the final feature maps; IG requires interpolation and multiple forward/backward passes (50 steps), making Grad-CAM faster for real-time clinical use. (3) \textbf{Transparency}: By explicitly leveraging the model's learned feature maps, Grad-CAM's reasoning is anchored to the architecture; IG's gradient-based signal is more abstract. (4) \textbf{Validation strategy}: With samples showing clear color-difference heatmaps, a clinician can immediately verify the model's focus; with those showing black differences, we flag them as unreliable. This creates an actionable decision rule: deploy Grad-CAM with internal reliability flagging, treating colored-difference samples as higher-confidence AI-assisted diagnoses. A complementary approach would be to retrain with architectural choices designed to be more faithful (e.g., attention layers, inherently interpretable models) and re-run the sanity checks, but given current tools, Grad-CAM offers the best balance of transparency, speed, and clinical applicability.
